\documentclass[10pt]{article}
\usepackage[T1]{fontenc}
\usepackage{mathpazo}
\renewcommand{\ttdefault}{pcr}
\sloppy
\usepackage[a4paper,margin=18mm]{geometry}
\usepackage{booktabs}
\usepackage{rotating}
\usepackage{amsmath}
\usepackage{graphicx}
\usepackage[hidelinks]{hyperref}
\usepackage{url}
\renewcommand{\thetable}{S\arabic{table}}
\renewcommand{\thesection}{S\arabic{section}}
\setlength{\parskip}{4pt}\setlength{\parindent}{0pt}

\title{Supplemental Material for\\ ``Container Overhead or Stale Dependencies? Decomposing the Latency and Energy Cost of Containerised Inference on an Edge Node''}
\author{Lin Ding Shan\\ UCSI University, Kuala Lumpur, Malaysia}
\date{}

\begin{document}
\maketitle

This document supports the main paper. Section numbers, equations and tables cited as ``paper'' refer to the main text. Every table here is generated from the raw run logs by \texttt{analysis/make\_supp.py} in the released repository; nothing is typed by hand.

\section{Run schedule, integrity checks and deviations}

Table~\ref{tab:s-runs} lists the nine runs in block order, with the within-block order as randomised (seed 20260927). Before each run, the orchestrating script (\texttt{orchestrate.py}) (i)~waited until the battery charge current was at most +50\,mA (``floating''); (ii)~waited until the SoC temperature was at most 50\,$^\circ$C; (iii)~from \texttt{legacy\_2} onward, disabled the camera's dynamic frame rate on the host with \texttt{v4l2-ctl} and saved the full camera control state next to the run; (iv)~launched the run; and afterwards (v)~checked that \texttt{RUN\_INFO.txt} reported completion and that the frame count lay within 0.93--1.05 of 129{,}600. A run failing any check stopped the orchestrator.

\textbf{Camera frame rate.} The USB camera runs in aperture-priority auto exposure. With \texttt{exposure\_dynamic\_framerate} enabled, it lowers its frame rate when the scene is dark. \texttt{legacy\_1} (03:56--06:56) and \texttt{matched\_1} (06:57--09:57) ran in a dark room; the median per-second rate was 15.0\,frame/s in every hour of both runs, and the overall rates were 15.046 and 15.049\,frame/s. The original \texttt{native\_1} run began at 09:58 (\texttt{orchestrator.log}); a check of its partial log during the run showed the frame rate rising towards 30\,frame/s after the room lights came on; the run was stopped, the orchestrator removed its directory before the replacement run, so this observation cannot be checked from the deposited data. All later runs, starting with \texttt{legacy\_2} and including the replacement \texttt{native\_1}, requested 15\,frame/s with dynamic frame rate disabled, and ran at 14.88--14.91\,frame/s. The loop is camera-bound in all nine runs, so the rate fixes the number of inferences per active phase.

\textbf{Harness update and image rebuild.} The harness used for \texttt{legacy\_1} and \texttt{matched\_1} did not request a frame rate and did not write the ``Measured rate'' line to \texttt{RUN\_INFO.txt}. It was updated before \texttt{legacy\_2} to add the \texttt{--fps} option and that line; the preprocessing, the timed \texttt{InferenceSession.run} call and the per-frame logging were not changed. The earlier version of the harness was not retained. The images copy the harness in their last layer, after the layers that install the stack, and were rebuilt to include it; unless the build cache is bypassed, such a rebuild reuses the stack layers. Image layer digests were not recorded. The environment fingerprints (Table~\ref{tab:s-fp}) were taken on 27 September at 02:03, before the first run, and were not repeated after the rebuild; the per-run \texttt{provenance.json} records show the same interpreter, glibc, ONNX Runtime, NumPy and OpenCV versions in all runs of each condition.

\textbf{Replicate field.} The ``Replicate'' line of each \texttt{RUN\_INFO.txt} holds the run's position in the nine-run schedule (1--9), not its replicate number within the condition; the directory label (e.g.\ \texttt{native\_2}) identifies the replicate.

\textbf{Block 1.} The replacement \texttt{native\_1} ran from 20:19 to 23:19 on the same day, after block~2 and immediately before block~3. Block~1 is therefore not a contiguous session. The paper's Section~4.5 repeats the analysis on blocks~2 and~3 only.

\textbf{Host state.} All systemd timers (package updates, manual-page indexing, swap write-back and others) were stopped for the whole campaign; the host software was not changed between fingerprinting and the last run. The graphical desktop session (\texttt{labwc} compositor, panel and file manager) remained running in every run and used under 1\% CPU. Provenance records show the same kernel (\texttt{6.18.39+rpt-rpi-2712}), governor (\texttt{ondemand}), four-CPU affinity and model SHA-256 (\texttt{bd79f604\ldots eb7ccb752}) in all nine runs; container runs recorded cgroup \texttt{cpu.max = max 100000} (no quota). The firmware throttle word, readable only on the host, was \texttt{0x0} before each native run; in all runs the 1\,Hz sysfs sampler, identical in all conditions, recorded no frequency cap and no under-voltage alarm (Table~\ref{tab:s-runs}).

\begin{table*}[!htbp]\centering\scriptsize
\caption{Run schedule and integrity. Start times are local (UTC+8). Frames: all inferences in the run; analysed: after the 600\,s warm-up. Battery current: mean (min, max) over the whole run, positive into the pack. Cap/UV: 1\,Hz samples with the CPU frequency capped below 2400\,MHz / under-voltage alarm set.}
\label{tab:s-runs}\setlength{\tabcolsep}{3.2pt}
\begin{tabular}{llllrrrrrrr}\toprule
Block & Run & Condition & Start & Frames & Analysed & Rate (frame/s) & Active / sleep (s) & Battery I (mA) & Bus V (V) & Cap / UV\\\midrule
1 & \texttt{legacy\_1} & legacy & 2026-09-27 03:56:12 & 130,093 & 122,878 & 15.046 & 60.05 / 15.02 & 5.0 (\textminus 12, 39) & 15.306 & 0 / 0\\
1 & \texttt{matched\_1} & matched & 2026-09-27 06:57:14 & 130,054 & 122,843 & 15.049 & 60.02 / 15.00 & 1.8 (\textminus 24, 27) & 15.296 & 0 / 0\\
1 & \texttt{native\_1} & native & 2026-09-27 20:19:29 & 128,922 & 121,775 & 14.910 & 60.06 / 15.00 & 0.9 (\textminus 27, 25) & 15.293 & 0 / 0\\
\addlinespace
2 & \texttt{legacy\_2} & legacy & 2026-09-27 11:15:22 & 128,865 & 121,720 & 14.906 & 60.04 / 15.02 & 1.1 (\textminus 21, 23) & 15.292 & 0 / 0\\
2 & \texttt{native\_2} & native & 2026-09-27 14:16:24 & 128,923 & 121,776 & 14.910 & 60.06 / 15.00 & 1.0 (\textminus 18, 25) & 15.298 & 0 / 0\\
2 & \texttt{matched\_2} & matched & 2026-09-27 17:17:27 & 128,895 & 121,748 & 14.907 & 60.05 / 15.00 & 1.2 (\textminus 17, 26) & 15.298 & 0 / 0\\
\addlinespace
3 & \texttt{matched\_3} & matched & 2026-09-27 23:22:31 & 128,788 & 121,639 & 14.896 & 60.04 / 15.00 & 0.7 (\textminus 27, 24) & 15.298 & 0 / 0\\
3 & \texttt{legacy\_3} & legacy & 2026-09-28 02:23:34 & 128,573 & 121,444 & 14.879 & 60.02 / 15.02 & 0.9 (\textminus 21, 24) & 15.298 & 0 / 0\\
3 & \texttt{native\_3} & native & 2026-09-28 05:24:36 & 128,577 & 121,448 & 14.877 & 60.03 / 15.00 & 1.0 (\textminus 20, 24) & 15.292 & 0 / 0\\
\addlinespace
\bottomrule\end{tabular}\end{table*}
\begin{sidewaystable*}[p]\centering\scriptsize\setlength{\tabcolsep}{1.9pt}
\caption{Per-run results after warm-up. Latency in ms; CPU ms: CPU time per inference; $P_\mathrm{floor}$: settled sleep-phase floor; energies per inference in mJ, Eqs.~(1)--(3) of the paper; temperatures in $^\circ$C (active-phase mean, mean cyclic peak, maximum); fan and frequency are active-phase means.}
\label{tab:s-perrun}
\begin{tabular}{lrrrrrrrrrrrrrrrrrrr}\toprule
Run & Lat. mean & median & p95 & p99 & SD & CPU \% & RAM \% & $P_\mathrm{act}$ W & $P_\mathrm{slp}$ W & $P_\mathrm{floor}$ W & $E_\mathrm{base}$ & $E_\mathrm{inc}$ & $E_\mathrm{cyc}$ & J/cycle & $T$ act. & $T$ peak & $T$ max & Fan rpm & MHz\\\midrule
\texttt{native\_1} & 19.51 & 20.30 & 26.30 & 28.90 & 3.85 & 56.24 & 8.59 & 10.541 & 7.020 & 6.674 & 559.5 & 265.2 & 824.6 & 738.4 & 64.27 & 66.54 & 68.3 & 5550 & 2399\\
\texttt{native\_2} & 19.24 & 20.00 & 26.40 & 29.70 & 4.01 & 55.93 & 8.58 & 10.453 & 6.972 & 6.617 & 554.7 & 263.3 & 817.9 & 732.4 & 61.45 & 63.68 & 65.5 & 5179 & 2399\\
\texttt{native\_3} & 19.38 & 20.30 & 25.30 & 28.10 & 3.77 & 56.07 & 8.60 & 10.381 & 6.911 & 6.554 & 550.7 & 263.2 & 813.9 & 726.8 & 62.14 & 64.43 & 65.5 & 5278 & 2399\\
\texttt{matched\_1} & 21.15 & 22.50 & 25.00 & 33.20 & 3.97 & 58.12 & 8.79 & 10.548 & 6.974 & 6.650 & 552.4 & 264.4 & 816.8 & 737.8 & 62.52 & 64.68 & 66.1 & 5309 & 2399\\
\texttt{matched\_2} & 20.50 & 20.00 & 26.80 & 30.50 & 4.15 & 57.01 & 8.79 & 10.500 & 6.929 & 6.564 & 550.4 & 270.1 & 820.5 & 734.5 & 61.95 & 64.22 & 67.2 & 5223 & 2398\\
\texttt{matched\_3} & 20.39 & 20.00 & 26.80 & 30.66 & 4.21 & 56.90 & 8.82 & 10.496 & 6.953 & 6.579 & 552.0 & 269.3 & 821.3 & 734.6 & 61.76 & 64.06 & 68.3 & 5171 & 2398\\
\texttt{legacy\_1} & 28.60 & 25.30 & 43.60 & 44.00 & 5.44 & 48.55 & 8.71 & 9.713 & 6.735 & 6.549 & 544.1 & 213.3 & 757.5 & 684.4 & 58.16 & 60.27 & 62.2 & 3907 & 2393\\
\texttt{legacy\_2} & 30.18 & 29.70 & 41.40 & 43.70 & 4.90 & 49.03 & 8.71 & 10.030 & 7.036 & 6.708 & 562.6 & 228.3 & 790.9 & 707.9 & 61.72 & 63.82 & 65.5 & 5302 & 2399\\
\texttt{legacy\_3} & 29.01 & 26.30 & 39.70 & 43.80 & 4.90 & 48.16 & 8.72 & 9.702 & 6.668 & 6.348 & 533.4 & 230.8 & 764.2 & 682.4 & 58.74 & 60.75 & 62.2 & 4354 & 2398\\
\bottomrule\end{tabular}\end{sidewaystable*}
\begin{table*}[!htbp]\centering\scriptsize
\caption{Within-block differences. Each effect in the paper is the mean of the three block differences shown; the 95\% interval uses $t_{0.975,2}=4.303$.}
\label{tab:s-blocks}\setlength{\tabcolsep}{3pt}
\begin{tabular}{l rrr r rrr r rrr r}\toprule
 & \multicolumn{4}{c}{$B-A$} & \multicolumn{4}{c}{$C-B$} & \multicolumn{4}{c}{$C-A$}\\
\cmidrule(lr){2-5}\cmidrule(lr){6-9}\cmidrule(lr){10-13}
Quantity & blk 1 & blk 2 & blk 3 & mean & blk 1 & blk 2 & blk 3 & mean & blk 1 & blk 2 & blk 3 & mean\\\midrule
Latency, mean (ms) & 1.637 & 1.262 & 1.013 & 1.304 & 7.449 & 9.678 & 8.619 & 8.582 & 9.086 & 10.940 & 9.632 & 9.886\\
Latency, p95 (ms) & \textminus 1.300 & 0.400 & 1.500 & 0.200 & 18.600 & 14.600 & 12.900 & 15.367 & 17.300 & 15.000 & 14.400 & 15.567\\
Latency, p99 (ms) & 4.300 & 0.800 & 2.562 & 2.554 & 10.800 & 13.200 & 13.138 & 12.379 & 15.100 & 14.000 & 15.700 & 14.933\\
CPU utilisation (pp) & 1.882 & 1.072 & 0.827 & 1.260 & \textminus 9.577 & \textminus 7.972 & \textminus 8.741 & \textminus 8.763 & \textminus 7.694 & \textminus 6.900 & \textminus 7.914 & \textminus 7.503\\
Power, active (W) & 0.007 & 0.047 & 0.115 & 0.056 & \textminus 0.835 & \textminus 0.470 & \textminus 0.794 & \textminus 0.700 & \textminus 0.828 & \textminus 0.423 & \textminus 0.679 & \textminus 0.643\\
Power, sleep (W) & \textminus 0.046 & \textminus 0.043 & 0.042 & \textminus 0.016 & \textminus 0.239 & 0.107 & \textminus 0.284 & \textminus 0.139 & \textminus 0.284 & 0.064 & \textminus 0.242 & \textminus 0.154\\
Energy/inf., above floor (mJ) & \textminus 0.75 & 6.84 & 6.03 & 4.04 & \textminus 51.07 & \textminus 41.77 & \textminus 38.44 & \textminus 43.76 & \textminus 51.82 & \textminus 34.93 & \textminus 32.42 & \textminus 39.72\\
CPU time/inference (ms) & 3.609 & 2.901 & 2.031 & 2.847 & \textminus 25.433 & \textminus 21.378 & \textminus 23.325 & \textminus 23.379 & \textminus 21.824 & \textminus 18.477 & \textminus 21.295 & \textminus 20.532\\
Energy/inf., cycle (mJ) & \textminus 7.85 & 2.56 & 7.37 & 0.69 & \textminus 59.30 & \textminus 29.56 & \textminus 57.05 & \textminus 48.64 & \textminus 67.15 & \textminus 27.00 & \textminus 49.69 & \textminus 47.95\\
Cyclic peak ($^\circ$C) & \textminus 1.854 & 0.532 & \textminus 0.373 & \textminus 0.565 & \textminus 4.409 & \textminus 0.393 & \textminus 3.308 & \textminus 2.703 & \textminus 6.263 & 0.140 & \textminus 3.681 & \textminus 3.268\\
Fan, active (rpm) & \textminus 240.9 & 43.2 & \textminus 106.4 & \textminus 101.4 & \textminus 1401.4 & 79.7 & \textminus 817.8 & \textminus 713.2 & \textminus 1642.3 & 122.9 & \textminus 924.2 & \textminus 814.5\\
\bottomrule\end{tabular}\end{table*}
\begin{table*}[!htbp]\centering\scriptsize
\caption{Environment fingerprints: SHA-256 of the compiled extension module imported by each numerical package, and the ONNX Runtime build string. Host and matched container are identical in every field.}
\label{tab:s-fp}
\begin{tabular}{lll}\toprule Package & Host = matched container & Legacy container\\\midrule
onnxruntime & 1.29.0 & 1.19.2\\
 & \texttt{5574a1e5d1dc5bcf1dff1e4ac2d41f8a} & \texttt{4e3e7d9ac54723d6602cf0258041a9ef}\\
 & \texttt{9cd7db5ad35661c28ba1433897ceb23e} & \texttt{79bb3c40eb28e9c63d4d105a7d1ce5bc}\\
numpy & 2.5.2 & 1.26.4\\
 & \texttt{72b6f54f37d34fdc2901fb770030bc9f} & \texttt{147aff3b8942204bbd523583c39b7e6c}\\
 & \texttt{c22f5b1e5f8d0eac9745ca5a349fecdb} & \texttt{5ddd20387afc9126e65183c0c53fc129}\\
opencv-python-headless & 5.0.0.93 & 4.10.0.84\\
 & \texttt{b69eef7f164dd20965ee3d88221d041c} & \texttt{2c06682a72c65e01542095239c3fbfc0}\\
 & \texttt{58d8f6a1e18c4c3bc29588286a7d6eee} & \texttt{80156803e5c80972ffad2e5f15424e23}\\
psutil & 7.2.2 & 6.1.0\\
 & \texttt{463d351dd4b4d93957d30a7b645dd0bc} & \texttt{9e0377b6a81b798b1c7b3a9388f3ef60}\\
 & \texttt{cd09276fd5d26c0ebd8c10e5088ba1a3} & \texttt{7ea9a1508c1ac4ff632c6b1f1e7aa192}\\
smbus2 (pure Python) & 0.6.1 & 0.5.0\\
Python / glibc & 3.13.5 / 2.41 & 3.9.23 / 2.31\\
\bottomrule\end{tabular}
\par\medskip\raggedright\scriptsize ONNX Runtime build string, host and matched: \path{git-branch=HEAD, git-commit-id=2e2543fbe9, fp8-kv-cache=1, build type=Release}\par
ONNX Runtime build string, legacy: \path{git-branch=HEAD, git-commit-id=ffceed9d44, build type=Release, cmake cxx flags: -DNDEBUG -Wp,-D_FORTIFY_SOURCE=2 -Wp,-D_GLIBCXX_ASSERTIONS -fstack-protector-strong -O3 -pipe -ffunction-sections -fdata-sections -Wno-restrict  -DCPUINFO_SUPPORTED}\par
\end{table*}


\section{Per-run results}
Table~\ref{tab:s-perrun} gives every run-level quantity used in the paper, and Table~\ref{tab:s-blocks} gives the within-block differences from which each effect and interval in the paper's Table~2 and Fig.~3 is computed. The effect for a quantity is the mean of the three block differences $d_1,d_2,d_3$, and its 95\% interval is $\bar d \pm t_{0.975,2}\, s_d/\sqrt3$.

\section{Environment fingerprints and image definitions}
Table~\ref{tab:s-fp} gives the fingerprints. The matched image definition (\texttt{Dockerfile.matched}) was generated by \texttt{make\_matched\_dockerfile.py} from the running host interpreter; it installs Debian's \texttt{python3} package in \texttt{debian:trixie-slim} and pins every package to the host's version. Fingerprints were taken with \texttt{fingerprint.py} on the host and inside both images, and the host and matched records compare equal as whole JSON documents.

\textbf{Legacy image.} The first campaign's image could not be recovered, so condition~C was rebuilt from its definition (\texttt{Dockerfile.legacy}). The original definition installed \texttt{onnxruntime} without a version constraint on \texttt{python:3.9-slim-bullseye}. On 64-bit ARM with CPython~3.9, the last ONNX Runtime release that publishes a compatible wheel is 1.19.2, so that version is determined. The other packages are pinned to the versions in the paper's Table~1; the first campaign's exact resolution of those packages cannot be established. The original definition also installed the system package \texttt{libglib2.0-0}, which the base distribution's archive no longer serves; it was omitted. The headless OpenCV build does not load it, and the harness imports and runs unchanged. The ONNX Runtime build strings (Table~\ref{tab:s-fp}) show that the host and legacy runtimes are different builds, with the legacy build recording its compiler flags.

\section{Instrumentation details}
\textbf{Per-frame log} (\texttt{basil\_data.csv}): wall-clock timestamp; latency of the \texttt{InferenceSession.run} call, measured with \texttt{time.time()} and logged at 0.1\,ms resolution; instantaneous frame rate; system-wide CPU utilisation since the previous frame (\texttt{psutil}); memory utilisation; SoC temperature; throttle flag; predicted class and confidence; the most recent 1\,Hz supply reading.

\textbf{1\,Hz log} (\texttt{power\_log.csv}), written through both phases: phase; Type-C bus voltage, current and power (UPS HAT (E) registers \texttt{0x10}, \texttt{0x12}, \texttt{0x14}); battery voltage, signed current and charge (\texttt{0x20}, \texttt{0x22}, \texttt{0x24}); SoC temperature; CPU utilisation over the previous second; current CPU frequency and frequency cap from \texttt{cpufreq}; the \texttt{rpi\_volt} low-voltage alarm; fan speed and PWM from \texttt{hwmon}. The sampler runs as a thread in the measured process and reads the same sysfs files in all three conditions.

\textbf{Cycle log} (\texttt{cycle\_events.csv}) marks each active and sleep transition. Phase durations in Table~\ref{tab:s-runs} are computed from it.

\textbf{Consistency check.} Run means of bus current and of bus power divided by bus voltage agreed within 0.4\% in every run. Individual samples disagree by more, because the module's registers are not updated together. The vendor does not specify the module's absolute accuracy.

\textbf{Sleep-phase floor.} The \texttt{ondemand} governor held the CPU at 2400\,MHz for about 5\,s after each active phase and then stepped down to 1500\,MHz; node power fell with it. The settled floor $P_\mathrm{floor}$ is the mean of sleep-phase samples taken at least 6\,s after the phase began. The energy split in the paper, Eqs.~(2)--(3), uses this floor. Table~\ref{tab:s-perrun} also gives the plain sleep-phase mean $P_\mathrm{slp}$.

\textbf{Warm-up and cycles.} The first 600\,s of every run, timed from the run's first cycle event, are excluded; the remaining 136 complete duty cycles per run are analysed. Cyclic peak temperature is the maximum 1\,Hz SoC temperature reading within each active phase (and the following 2\,s), averaged over the 136 cycles.

\section{Camera settings}
The saved control state (\path{runs/<label>.camera.txt}, the seven runs from \texttt{legacy\_2} onward) shows the same settings in every run:
\begin{itemize}\setlength{\itemsep}{0pt}
\item \path{auto_exposure = 3} (aperture priority);
\item \path{exposure_dynamic_framerate = 0};
\item \path{gain = 255}, \path{power_line_frequency = 2} (60\,Hz);
\item automatic white balance and continuous autofocus enabled;
\item 640$\times$480 at a requested 15\,frame/s.
\end{itemize}
The anti-flicker setting is 60\,Hz, whereas the local mains frequency is 50\,Hz. This may affect image exposure under artificial light, but it was the same in all runs and does not enter the timed inference call.

\section{Reproducing the analysis}
From the repository root, with the deposited runs extracted to \texttt{runs/}:
\begin{verbatim}
python3 analysis/analyse_v2.py     # -> out/numbers_v2.json
python3 analysis/make_tex.py       # -> out/numbers_v2.tex, out/table_results.tex
python3 analysis/make_supp.py      # -> out/supp_tables.tex
python3 analysis/figures_v2.py     # -> figures/fig1..fig5 (.pdf, .png)
\end{verbatim}
The paper's text, abstract and tables take their numbers from \texttt{numbers\_v2.tex} and \texttt{table\_results.tex}, so a rebuild from a fresh analysis cannot silently diverge from the data; the first-campaign values are read from the repository's earlier analysis (\path{analysis/numbers.json}, produced by \path{analysis/recompute.py}). Fixed parameters of the analysis (the 600\,s warm-up, the 6\,s settling time and the 17.5\,ms boundary between latency modes) are set in \texttt{analyse\_v2.py}. \texttt{manuscript/build.sh} runs the whole chain and compiles both documents. The run kit (\texttt{runner/}) contains the harness, orchestrator, image definitions and fingerprinting script as used in the campaign, and an as-executed description of the run procedure, including the changes made during the campaign (\texttt{PROTOCOL.md}).

\end{document}
